\chapter{Tag 14 — Encoder zusammenstecken}

\begin{quote}
In den letzten drei Etappen hast du jedes einzelne Teil eines
Datamatrix-Encoders studiert: ASCII-Codewords (Tag 11), Reed-Solomon-
Prüfbytes (Tag 12), Datenplatzierung im Symbol (Tag 13). Heute
verbindest du alle Theorie-Stücke zu einem Programm, das aus einem
Wort ein scannbares PNG erzeugt — und das Handy hält den
Daumen hoch.
\end{quote}

\section*{Lernziele für heute}

Am Ende von Tag 14 kannst du \dots

\begin{itemize}
\item den kompletten Encoder-Datenfluss in einem einzigen Skript
  hintereinanderschreiben (Stufe „roh") und Schritt für Schritt
  beobachten, wie aus dem Wort \texttt{GRETA} ein PNG wird
\item ASCII-Encoding und einen Reed-Solomon-Encoder inline
  implementieren — als kompakte Wiederholung der Theorie aus
  Tagen 11 und 12
\item das L-Pattern und das Zeitsignal selbst um den Datenbereich
  legen, der aus \texttt{place\_codewords} (Tag 13) kommt
\item ein Pillow-Bild mit korrekter Quiet-Zone speichern, das ein
  Smartphone-Scanner ohne Murren liest
\item den Bogen schließen: dein handgezeichneter \texttt{GRETA}-
  Datamatrix von gestern wird genauso gut gelesen wie das Code-
  generierte PNG von heute
\end{itemize}

\paragraph{Material} Laptop mit Thonny und installiertem
\texttt{pylibdmtx}, dein Smartphone mit einer Datamatrix-App
(zum Beispiel „QR \& Barcode Scanner"), dein
handgezeichnetes \texttt{GRETA}-Symbol vom Werkstattbogen aus
Etappe 13, ein dicker schwarzer Filzstift (falls Module nachgezogen
werden müssen).

% =============================================================
\section{Was du heute zusammensteckst \normalfont(\textit{≈ 5 min})}
% =============================================================

In Tagen 11--13 haben wir den Encoder in Einzelteilen gebaut:

\begin{center}
\begin{tikzpicture}[
  node distance=3mm and 6mm,
  font=\footnotesize,
  baustein/.style={draw, rounded corners=2pt, fill=gray!10, inner sep=4pt,
    minimum height=8mm, align=center, text width=24mm},
  pfeil/.style={->, >=stealth, thick},
]
  \node[baustein] (text)  {Klartext\\\texttt{"GRETA"}};
  \node[baustein, right=of text]  (ascii) {Tag 11\\ASCII +1};
  \node[baustein, right=of ascii] (rs)    {Tag 12\\Reed-Solomon};
  \node[baustein, right=of rs]    (place) {Tag 13\\Platzierung};
  \node[baustein, right=of place] (frame) {\textbf{heute}\\Frame + PNG};

  \draw[pfeil] (text)  -- (ascii);
  \draw[pfeil] (ascii) -- (rs);
  \draw[pfeil] (rs)    -- (place);
  \draw[pfeil] (place) -- (frame);

  \node[below=1mm of ascii, font=\scriptsize, gray] {5 Datenbytes};
  \node[below=1mm of rs,    font=\scriptsize, gray] {+ 7 ECC = 12};
  \node[below=1mm of place, font=\scriptsize, gray] {Datenmodule};
  \node[below=1mm of frame, font=\scriptsize, gray] {scannbar};
\end{tikzpicture}
\end{center}

Heute liest du nicht über jeden Schritt nochmal in der Tiefe —
stattdessen schreibst du sie in einem einzigen Skript hintereinander
und siehst zu, wie der Datenfluss durchläuft. Das ist die „rohe"
Stufe: alles auf Top-Level, kein Schnickschnack, kein Refactoring.
Sie hat den Charme, dass du das ganze Mosaik in einem Stück siehst.
In Etappe 15 räumen wir mit Funktionen und einer Klasse auf.

% =============================================================
\section{Schritt 1 — Klartext zu Datenbytes \normalfont(\textit{≈ 5 min})}
% =============================================================

Datamatrix kennt mehrere Encoding-Modi (ASCII, C40, Text, Base 256,
\dots). Wir bleiben beim einfachsten — ASCII-Modus, ein Byte pro
Zeichen mit Offset $+1$:

\begin{pythoneinheit}{1}{ASCII-Modus inline}
\pythoncodeteil{tag14/encoder_roh.py}{ascii}

Für \texttt{GRETA} ergibt das die fünf Bytes
\texttt{[72, 83, 70, 85, 66]}. Das sind genau die fünf Datenbytes,
die ein 12$\times$12-Symbol aufnimmt — wir brauchen also nicht zu
padden.
\end{pythoneinheit}

\begin{bleistiftuebung}{1}{ASCII von Hand}
Rechne die Bytes für dein eigenes 5-Buchstaben-Wort aus (groß­geschrieben,
ohne Umlaute). Notiere sie auf dem Karopapier — du brauchst sie
gleich, um deinen Encoder mit deinem eigenen Wort laufen zu lassen.
\end{bleistiftuebung}

% =============================================================
\section{Schritt 2 — Reed-Solomon-ECC \normalfont(\textit{≈ 15 min})}
% =============================================================

Aus den 5 Datenbytes berechnen wir 7 Prüfbytes (12$\times$12-Symbole
verwenden immer 7 ECC-Codewords). Wir nehmen nicht den fertigen
Encoder aus Tag 12, sondern bauen das Wesentliche \emph{nochmal}
inline — als Wiederholung und als Beweis, dass es kein Hexenwerk ist.

\begin{pythoneinheit}{2}{GF(2$^8$)-Multiplikation und Generator-Polynom}
\pythoncodeteil{tag14/encoder_roh.py}{ecc}

Der Trick im letzten Block ist die Polynomdivision: wir hängen 7
Nullen an die Datenfolge an und ziehen Schritt für Schritt das
Generator-Polynom mal dem aktuellen Leitkoeffizienten ab. Was am
Ende übrig bleibt, sind die 7 ECC-Bytes.
\end{pythoneinheit}

\paragraph{Erwartetes Ergebnis} Wenn du das Skript laufen lässt,
sollten die ECC-Bytes für \texttt{GRETA} so aussehen:

\begin{center}
\texttt{ECC = [64, 90, 71, 128, 186, 106, 125]}
\end{center}

Das sind exakt die Werte, die du gestern in deine Hilfstabelle auf
dem Werkstattbogen eingetragen hast (Tag 13, Lösung zu
Bleistift-Übung 3). Wenn die Zahlen abweichen, liegt der Fehler
fast immer beim Generator-Polynom-Aufbau oder an der Reihenfolge
der Koeffizienten in der Polynomdivision.

% =============================================================
\section{Schritt 3 — Datenmodule platzieren \normalfont(\textit{≈ 5 min})}
% =============================================================

Den Platzierungs-Algorithmus haben wir gestern ausführlich in
\texttt{place\_codewords} implementiert (Tag 13). Heute rufen wir ihn
einfach auf:

\begin{pythoneinheit}{3}{Datenraster aus den 12 Codewords}
\pythoncodeteil{tag14/encoder_roh.py}{matrix}

Die zurückgegebene 12$\times$12-Matrix enthält jetzt 0/1-Werte für
alle Datenmodule und \texttt{None} dort, wo der Frame hinkommt
(L-Pattern und Zeitsignal — das ist Schritt 4).
\end{pythoneinheit}

% =============================================================
\section{Schritt 4 — L-Pattern und Zeitsignal \normalfont(\textit{≈ 10 min})}
% =============================================================

Auf dem Werkstattbogen gestern war der Frame schon vorgedruckt; in
einem reinen Code-Encoder müssen wir ihn selbst setzen.

\begin{itemize}
\item \textbf{L-Pattern}: untere Kante (Zeile~11) und linke Kante
  (Spalte~0) komplett schwarz. Daran erkennt der Scanner, wo das
  Symbol anfängt und wie groß es ist.
\item \textbf{Zeitsignal}: obere Kante alterniert (gerade Spalten
  schwarz, ungerade weiß), rechte Kante alterniert (ungerade Zeilen
  schwarz, gerade weiß). Daran zählt der Scanner die Module ab.
\end{itemize}

\begin{bleistiftuebung}{2}{Frame ohne Code zeichnen}
Zeichne auf einem leeren 12$\times$12-Karoraster nur den Frame
(L-Pattern und Zeitsignal), ohne irgendeinen Datenmodul. Das
Ergebnis sollte aussehen wie ein leeres Datamatrix-„Skelett".
Vergleiche mit der Vorlage von gestern — die Frame-Module sind
identisch, weil sie nicht vom Inhalt abhängen.
\end{bleistiftuebung}

\begin{pythoneinheit}{4}{Frame in die Matrix eintragen}
\pythoncodeteil{tag14/encoder_roh.py}{frame}

Vier kompakte Schleifen — keine Magie. Was vorher \texttt{None} war,
ist jetzt 0 oder 1.
\end{pythoneinheit}

% =============================================================
\section{Schritt 5 — Aus der Matrix wird ein PNG \normalfont(\textit{≈ 10 min})}
% =============================================================

Pillow kennst du schon aus Tag 10 (EAN-13). Jetzt das gleiche Spiel
für Datamatrix:

\begin{pythoneinheit}{5}{Pillow: 12$\times$12-Pixel und skaliertes Druck-PNG}
\pythoncodeteil{tag14/encoder_roh.py}{render}

Zwei wichtige Details fürs Scannen:

\begin{enumerate}
\item \textbf{Quiet-Zone} — der weiße Rand drum­herum. Ohne den
  weiß der Scanner nicht, wo das Symbol aufhört. Faustregel:
  mindestens eine Modulbreite Rand auf jeder Seite.
\item \textbf{Modus „L"} statt „1" — die meisten Scanner-Bibliotheken
  (auch \texttt{pylibdmtx}) lehnen reine 1-Bit-PNGs ab. „L" ist
  8-Bit-Graustufen, gleicher visueller Inhalt.
\end{enumerate}
\end{pythoneinheit}

% =============================================================
\section{Scan-Test \normalfont(\textit{≈ 15 min})}
% =============================================================

Erst die Selbstkontrolle mit \texttt{pylibdmtx}, dann der echte
Scan mit dem Handy.

\begin{pythoneinheit}{6}{Selbstkontrolle: pylibdmtx liest das PNG}
\begin{minted}{python}
from pylibdmtx.pylibdmtx import decode
from PIL import Image

img = Image.open("greta.png")
print(decode(img))
\end{minted}

Erwartet:
\begin{minted}{text}
[Decoded(data=b'GRETA', rect=Rect(left=29, top=30, width=300, height=299))]
\end{minted}
\end{pythoneinheit}

\begin{bleistiftuebung}{3}{Smartphone-Scan: Code vs. Hand}
Drucke deine Datei \texttt{greta.png} aus (oder zeige sie auf dem
Bildschirm) und scanne sie mit der Datamatrix-App auf deinem Handy.
Mache dann das gleiche mit deinem \emph{handgezeichneten}
\texttt{GRETA}-Symbol vom Werkstattbogen aus Tag 13. Bei beiden
sollte „GRETA" zurückkommen.

Falls die Hand-Zeichnung nicht erkannt wird, hilft es meist, sie
abzufotografieren und das Foto am Handy nochmal zu scannen — dann
korrigiert die Kamera kleine Schiefstellungen automatisch. Notfalls
einzelne Module mit dem dicken Filzstift nachziehen.
\end{bleistiftuebung}

% =============================================================
\section{Aufräumen mit Funktionen \normalfont(\textit{≈ 20 min})}
% =============================================================

Stufe roh war prima zum Verstehen — aber unhandlich, sobald du
mehr willst als ein einziges fest verdrahtetes Wort. Probleme:

\begin{itemize}
\item Mit anderem Text encoden? Konstante \texttt{TEXT} oben
  ändern und das Skript neu starten.
\item Mehrere Wörter in einer Schleife encoden? Geht nicht ohne
  copy-pasten.
\item Modulgröße fürs Drucken anpassen? Im Render-Block
  manuell die \texttt{300} und \texttt{30} und \texttt{360}
  ändern und hoffen, dass alles zueinander passt.
\end{itemize}

Die naheliegende Antwort: jeder logische Schritt wird eine eigene
Funktion. Datei \texttt{encoder\_funktionen.py}.

\begin{pythoneinheit}{7}{Schritt 1 + 2 als Funktionen}
\pythoncodeteil{tag14/encoder_funktionen.py}{ascii}
\pythoncodeteil{tag14/encoder_funktionen.py}{rs}

Die GF(2$^8$)-Multiplikation und der Generator-Polynom-Aufbau
liegen jetzt in eigenen Helper-Funktionen — wir müssen sie nicht
wieder hineinschreiben, wenn wir eine zweite Encoder-Variante
basteln.
\end{pythoneinheit}

\begin{pythoneinheit}{8}{Matrix bauen + PNG schreiben als Funktionen}
\pythoncodeteil{tag14/encoder_funktionen.py}{matrix}
\pythoncodeteil{tag14/encoder_funktionen.py}{render}

\texttt{matrix\_zu\_png} bekommt jetzt zwei Stellschrauben:
\texttt{modulgroesse} (wie viele Pixel ein Modul groß ist) und
\texttt{quietzone\_module} (wie viele Module weißer Rand). Damit
kannst du das gleiche Symbol für Bildschirm-Vorschau klein und
für Druck groß rendern, ohne den Code anfassen zu müssen.
\end{pythoneinheit}

\begin{pythoneinheit}{9}{Pipeline + Demo: drei Wörter in einer Schleife}
\pythoncodeteil{tag14/encoder_funktionen.py}{pipeline}
\pythoncodeteil{tag14/encoder_funktionen.py}{demo}

Der Demo-Block bringt drei Wörter — \texttt{GRETA}, \texttt{HONIG},
\texttt{BIENE} — in einem Rutsch zu drei PNG-Dateien. In Stufe roh
hättest du dafür dreimal das ganze Skript laufen lassen müssen,
jedes Mal die Konstante \texttt{TEXT} angepasst.
\end{pythoneinheit}

\begin{bleistiftuebung}{4}{Eigenes Test-Set}
Lass das Skript laufen und prüfe alle drei PNGs am Handy. Ergänze
dann die Wortliste um dein eigenes 5-Buchstaben-Wort und um eines
deiner Familienmitglieder (jeweils \emph{exakt} fünf Großbuchstaben,
ohne Umlaute). Du solltest danach fünf scanbare Symbole im
Verzeichnis liegen haben.
\end{bleistiftuebung}

% =============================================================
\section{Reflexion und Brücke zu Tag 15 \normalfont(\textit{≈ 10 min})}
% =============================================================

Du hast jetzt zwei Versionen desselben Encoders gebaut:
\texttt{encoder\_roh.py} (ein Skript, alles top-level) und
\texttt{encoder\_funktionen.py} (sauber in Funktionen aufgeteilt).
Die zweite Version kann mehr — Schleifen, einstellbare Modulgröße,
Wiederverwendung der GF-Helper.

\subsection*{Was an der Funktions-Version trotzdem noch fehlt}

\begin{itemize}
\item Wir können nur 12$\times$12. Andere Symbol-Größen
  (14$\times$14, 16$\times$16, 18$\times$18, \dots) brauchen
  andere Codeword-Anzahlen, mehr ECC-Blöcke, andere
  Padding-Regeln. Der jetzige Code hat \texttt{anzahl\_pruefbytes=7}
  und \texttt{n=12} fest verdrahtet.
\item Wenn jemand zwei Encoder-Varianten parallel benutzen
  möchte (z. B. einen für 12$\times$12 und einen für
  14$\times$14), wird er pro Aufruf alle Parameter neu durchreichen
  müssen — eine Klasse mit Konfiguration im \texttt{\_\_init\_\_}
  wäre handlicher.
\item Encode\,$\to$\,RS\,$\to$\,Matrix\,$\to$\,PNG ist immer der
  gleiche Ablauf, aber jeder Aufrufer muss ihn wieder
  hinschreiben (oder \texttt{datamatrix\_encoden} aufrufen, das
  diese Reihenfolge versteckt — ein erster Schritt Richtung
  Klasse).
\end{itemize}

\subsection*{Vorschau Tag 15}

Wir machen aus den Funktionen eine \texttt{DatamatrixEncoder}-Klasse
— sauber gekapselt mit Methoden für jeden Schritt. Damit wird der
Encoder zum Werkzeug:
\texttt{DatamatrixEncoder("BLAUBEERE").save("blaubeere.png")}.
Und dann erweitern wir auf größere Symbole, sodass auch längere
Texte hineinpassen.
